\chapter{线程、锁、原子操作与内存模型}

推理引擎要使用多核，就必须处理并发。并发不是简单地把循环丢给多个线程。线程之间会共享模型、共享输入、写输出、更新队列、等待任务完成。只要存在共享状态，就会遇到锁、原子操作、内存可见性和数据竞争问题。性能工程里，错误的并发代码比慢代码更危险，因为它可能偶发出错，并且很难复现。

\section{线程解决什么问题}

线程让多个执行流在同一进程地址空间内运行。对推理引擎而言，线程常用于并行计算输出通道、并行处理 batch、并行处理 attention head 或并行执行独立算子。线程共享模型权重可以避免复制大对象，但共享也意味着必须明确写入所有权。

第一阶段 LCQI 的线性层可以按输出神经元切分。每个线程负责一段输出下标，读取同一份输入和权重，写不同输出位置。这种切分简单、安全，适合教学。按输入维度切分则需要多个线程写同一个输出的部分和，最后再归约，复杂度更高。

\section{锁和临界区}

锁保护临界区，保证同一时刻只有一个线程修改共享状态。任务队列、日志、全局统计、内存池元数据都可能需要锁。锁的成本不只在加锁和解锁本身，还包括线程阻塞、唤醒、缓存行争用和优先级反转。临界区越小越好，共享写越少越好。

\begin{lstlisting}[language=C++,caption={用互斥锁保护共享计数器}]
#include <cstdint>
#include <mutex>

class SafeCounter {
public:
    void add(std::int64_t value) {
        std::lock_guard<std::mutex> guard(this->mutex_);
        this->value_ += value;
    }

    std::int64_t value() const {
        std::lock_guard<std::mutex> guard(this->mutex_);
        return this->value_;
    }

private:
    mutable std::mutex mutex_;
    std::int64_t value_ = 0;
};
\end{lstlisting}

这个例子使用了明确的 \code{this->} 成员访问，适合小范围共享状态。真正的热路径不应频繁进入这种临界区。如果每个输出元素都加锁一次，推理引擎会被同步成本拖垮。正确做法通常是线程本地累计，最后合并。

\section{原子操作和内存序}

原子操作可以在不使用互斥锁的情况下安全读写单个共享变量。它适合计数器、状态标志、任务索引等简单场景。原子操作仍然有成本，并且内存序会影响可见性和重排。默认的 sequentially consistent 最容易理解，但可能比必要约束更强；relaxed 适合只需要原子性、不需要同步其他内存的计数器。

\begin{lstlisting}[language=C++,caption={任务索引用原子递增分发}]
#include <atomic>
#include <cstdint>

class AtomicTaskCursor {
public:
    explicit AtomicTaskCursor(std::int32_t end) : end_(end) {}

    bool next(std::int32_t* task) {
        const std::int32_t current =
            this->cursor_.fetch_add(1, std::memory_order_relaxed);
        if (current >= this->end_) {
            return false;
        }
        *task = current;
        return true;
    }

private:
    std::atomic<std::int32_t> cursor_{0};
    std::int32_t end_ = 0;
};
\end{lstlisting}

这里使用 relaxed 是因为任务编号本身不发布其他数据，所有任务数据在启动线程前已经准备好。若一个线程写入数据并通过原子标志通知另一个线程读取，就需要 release/acquire 语义。内存模型的核心问题是：一个线程写的内容，另一个线程何时保证看见。

\section{线程池和任务切分}

推理引擎不能每次算子调用都创建线程。线程池提前创建工作线程，算子提交任务，线程池执行并等待完成。线程池设计要控制任务粒度，避免过度切分；要避免所有线程抢同一个锁；要明确异常和停止机制。第一阶段可以实现简单 parallel for，后续再引入 work stealing 或 NUMA 感知。

\begin{keyidea}
并发优化的第一原则是写入所有权清晰。让每个线程写独立输出块，通常比让多个线程共同更新同一结果更容易正确，也更容易获得性能。
\end{keyidea}

\section{推理引擎里的并发边界}

LCQI 中，模型权重加载后只读，可以在线程间共享；输入张量一次推理内只读；输出张量按区间分配；任务队列和 benchmark 统计需要同步；日志不应出现在热路径。这个边界要在代码接口里体现出来，否则后续算子越多，数据竞争越难排查。

\begin{exercisebox}
把 LCQI 的隐藏层输出按区间切给多个线程。要求每个线程只写自己的输出区间，主线程等待所有任务完成后再进入下一层。比较 1、2、4 个线程的时间，并解释加速或不加速的原因。
\end{exercisebox}
